% ============================================================================
% Fichier  : partie5/P05_C13_sous_traitants.tex
% Titre    : Sous-traitants et Responsables Conjoints
% Auteur   : MINKA MI NGUIDJOÏ Thierry Emmanuel
% Version  : 1.0 -- novembre 2026
% Corrigés : corriges/Corr_C13.tex
% ============================================================================
% Sources : loi n° 2024/017, art. 5, 16, 20, 23(3), 24, 26, 30, 31 (texte
% officiel) ; ISO/IEC 27701 et 27018 (paraphrasées).
% Faits nouveaux (bible) : contrat SOREC (2019), contrat de l'éditeur du
% logiciel bancaire, avenant MEDISOFT du 20/05/2026, convention avec SAWA
% ASSURANCES pour l'assurance des emprunteurs ; 17 au 21 août 2026.
% ============================================================================

\chapter{Sous-traitants et Responsables Conjoints}
\label{chap:soustraitants}
\chaptermark{Sous-traitants et responsables conjoints}

\epigraph{%
    « Confier ses données à un prestataire,\\
    ce n'est pas lui confier sa responsabilité. »%
}{Principe du cours}

\niveauChapitre{Intermédiaire}{%
    Chapitres~\ref{chap:donnees} et~\ref{chap:cartographie}
}

\begin{boxobjectifs}
\begin{itemize}
    \item Appliquer les obligations des articles~16, 26 et~30 à un contrat
          de sous-traitance.
    \item Contrôler les clauses d'un contrat contre une liste de vérification
          déclarée.
    \item Vérifier que le sous-traitant applique ce que le contrat prévoit.
    \item Reconnaître une situation de responsables conjoints et vérifier la
          convention de l'article~31.
    \item Tirer les conséquences de l'article~20~: mêmes obligations pour le
          responsable et le sous-traitant.
\end{itemize}
\end{boxobjectifs}

\bigskip

% ============================================================================
\section*{Avant de commencer : un contrat de deux pages}
% ============================================================================

Douala, lundi 17~août~2026. Mme~BELLO rencontre le gérant de SOREC, le
prestataire de recouvrement de NYANGA CRÉDIT. Chaque mois, SOREC reçoit par
messagerie un classeur des débiteurs en retard~: noms, téléphones,
adresses, employeurs, montants. Ses agents appellent, se déplacent,
parfois chez l'employeur.

Le contrat, signé en 2019, tient sur deux pages~: objet, rémunération au
pourcentage des sommes recouvrées, durée. Il ne dit rien des données~:
ni leur usage, ni leur sécurité, ni leur sort à la fin du dossier. Le gérant
reconnaît sans difficulté que les classeurs des cinq dernières années sont
conservés « au cas où », et que ses agents ont accès à tous les dossiers de
tous les clients de SOREC.

\begin{boxalert}
\textbf{Violation typique.} Tenir le prestataire pour seul responsable de ce
qu'il fait des données, ou tenir le contrat pour une preuve que le
prestataire protège les données.

\smallskip
\textbf{Règle~:} le responsable « veille au respect » des mesures par son
sous-traitant (article~16) et doit s'assurer qu'il offre des garanties
suffisantes (article~30). Le contrat fixe les obligations~; l'audit vérifie
qu'elles sont tenues.
\end{boxalert}

% ============================================================================
\section{Les Obligations de la Loi}
\label{sec:c10-obligations}
% ============================================================================

\begin{boxjuris}
\textbf{\loicam{n°~2024/017}, article~30}~\cite{loi2024017}~:
\textit{« (1) Lorsque le responsable du traitement recourt à un
sous-traitant, il s'assure que celui-ci dispose des garanties suffisantes
quant à la mise en œuvre des mesures techniques et organisationnelles
appropriées de manière à ce que le traitement respecte les exigences de la
présente loi. (2) Le traitement par un sous-traitant est régi par un
contrat qui lie le sous-traitant au responsable du traitement. (3) Le
contrat visé à l'alinéa 2 ci-dessus définit, notamment, le type de données
à caractère personnel et les catégories de personnes concernées, les
obligations et les droits du responsable du traitement, ainsi que du
sous-traitant. »}

\smallskip
\textbf{Portée.} Le contrat est obligatoire~; son contenu minimal est
court. Le mot « notamment » laisse au responsable le soin de le compléter
pour satisfaire l'alinéa~1.
\end{boxjuris}

Trois autres articles complètent le régime. L'article~16 oblige le
sous-traitant à « apporter des garanties suffisantes au respect des
mesures de sécurité » et le responsable à veiller à leur application~;
l'article~26 répète l'obligation de garanties pour la sécurité technique et
organisationnelle~; l'article~20 assujettit le responsable et le
sous-traitant « aux mêmes obligations ». L'article~23, alinéa~3, les rend
enfin « solidairement responsables de toute publication de données à
caractère personnel effectuée en l'absence du consentement de la personne
concernée ».

Le sous-traitant doit aussi tenir son propre registre (article~29) et
informer l'Autorité et les personnes d'une violation (article~22). Un
prestataire qui ignore ces obligations n'est pas seulement un risque pour
son client~: il est lui-même en infraction.

% ============================================================================
\section{Le Contrat de Sous-traitance}
\label{sec:c10-contrat}
% ============================================================================

La loi ne fixe que trois éléments du contrat. Pour satisfaire l'obligation
de garanties suffisantes, le cours retient une liste de vérification de
douze clauses, déclarée comme critère d'audit et inspirée des normes
ISO/IEC~27701 et~27018.

\begin{tableaucours}{cL{5.2cm}L{5.8cm}}{%
    \tableauHeader{N°} &
    \tableauHeader{Clause} &
    \tableauHeader{Fondement}
}
1 & Types de données et catégories de personnes & Article~30(3) \\
2 & Obligations et droits des deux parties & Article~30(3) \\
3 & Objet, finalité et durée du traitement & Article~11 \\
4 & Traitement sur instructions documentées seulement & Article~24(2) \\
5 & Confidentialité des personnes autorisées & Articles~24 et~25(3) \\
6 & Mesures de sécurité, renvoi à l'article~27 & Articles~16, 26, 27 \\
7 & Sous-traitance ultérieure soumise à accord & Article~30(1) \\
8 & Assistance pour les droits des personnes & Articles~37 à~47 \\
9 & Information sans délai en cas de violation & Article~22 \\
10 & Restitution ou destruction en fin de contrat & Article~13 \\
11 & Droit d'audit du responsable & Article~16(2) \\
12 & Lieu du traitement, transferts éventuels & Article~32 \\
\end{tableaucours}
\vspace{-0.8em}
\begin{center}{\small\itshape Tableau~13.1. Les douze clauses du cours (critère déclaré)}\end{center}

Les clauses~1 et~2 sont exigées par la lettre de la loi~: leur absence est
un \NON{} sans discussion. Les dix autres traduisent l'obligation de
garanties~: leur absence est un \NON{} au regard du critère déclaré, et le
rapport le précise. Le laboratoire APDC-LAB-06 exerce cette lecture sur
cinq contrats.

\begin{boxinfo}
\textbf{Le droit d'audit.} La clause~11 est la condition de toutes les
autres. Sans droit d'audit, le responsable ne peut pas vérifier les
garanties qu'il doit s'assurer d'obtenir~; il ne peut que croire le
prestataire. Un prestataire qui refuse tout audit oblige l'auditeur à
coter \NEK{} chaque exigence qui dépend de lui, et le responsable à en
tirer les conséquences.
\end{boxinfo}

% ============================================================================
\section{Vérifier l'Exécution}
\label{sec:c10-execution}
% ============================================================================

Un contrat parfait ne prouve rien de ce que fait le prestataire. L'audit des
sous-traitants procède en trois cercles, du moins coûteux au plus probant~:

\begin{enumerate}
    \item \textbf{Les documents du prestataire}~: registre, politique de
          sécurité, certificats, rapports d'audit indépendants. Un
          certificat ISO/IEC~27001 ou 27018 couvre un périmètre~; on vérifie
          qu'il inclut le service fourni.
    \item \textbf{Le questionnaire et l'entretien}, recoupés par des pièces~:
          liste des personnes ayant accès, journal des accès aux données du
          client, preuve des destructions.
    \item \textbf{L'audit sur place}, exercé au titre de la clause~11, pour
          les prestataires qui portent les traitements les plus risqués.
\end{enumerate}

Le choix du cercle suit le risque. Un fournisseur de papeterie qui ne
reçoit qu'une adresse de livraison ne justifie qu'un contrôle documentaire~;
un hébergeur de dossiers médicaux justifie un audit sur place.

% ============================================================================
\section{Les Responsables Conjoints}
\label{sec:c10-conjoints}
% ============================================================================

\begin{boxjuris}
\textbf{\loicam{n°~2024/017}, article~31}~\cite{loi2024017}~:
\textit{« (1) En cas de traitement conjoint, les responsables déterminent
les finalités, ainsi que les moyens du traitement et en sont solidairement
responsables. (2) Les responsables conjoints du traitement définissent de
manière transparente, par une convention rendue publique, leurs
obligations respectives aux fins d'assurer le respect des exigences de la
présente loi, notamment en ce qui concerne leurs rôles respectifs, leurs
relations vis-à-vis de la personne concernée, l'exercice des droits de la
personne concernée, ainsi que leurs obligations respectives quant à la
communication des informations visées à l'article 21 ci-dessus.
(3) Nonobstant les termes de la convention visée à l'alinéa 2 ci-dessus,
la personne concernée exerce les droits que lui confère la présente loi à
l'égard des responsables du traitement. »}

\smallskip
\textbf{Portée.} La convention doit être \textit{rendue publique}~: c'est
une exigence plus forte que celle du droit européen, qui n'impose que de
mettre « les grandes lignes » à disposition. La personne peut s'adresser à
l'un ou l'autre responsable, quelle que soit la répartition convenue.
\end{boxjuris}

Il y a responsabilité conjointe quand deux organisations décident ensemble
de la finalité et des moyens, et pas quand l'une exécute les instructions de
l'autre. NYANGA CRÉDIT et SAWA ASSURANCES\footnote{SAWA ASSURANCES est
fictive.} proposent ensemble une assurance des emprunteurs~: NYANGA
collecte le questionnaire de santé à l'agence, SAWA ASSURANCES décide de
l'acceptation du risque, et les deux ont conçu ensemble le formulaire et
le circuit. Ce sont des responsables conjoints, et le questionnaire de
santé est une donnée sensible.

% ============================================================================
\section{Les Contrats de NYANGA et de MBOLO}
\label{sec:c10-cas}
% ============================================================================

\begin{boxscenario}{Missions NYANGA et MBOLO. Cinq contrats, 17 au 21 août 2026}
\begin{tableaucours}{L{3.4cm}cccL{4.2cm}}{%
    \tableauHeader{Contrat} &
    \tableauHeader{Cl. 1--2} &
    \tableauHeader{Cl. 3--12} &
    \tableauHeader{Audit} &
    \tableauHeader{Observation}
}
SOREC (recouvrement) & 0/2 & 0/10 & Non & Classeurs de cinq ans
conservés \\
Éditeur du logiciel bancaire & 2/2 & 6/10 & Non & Hébergement à Abidjan,
pas de clause de transfert \\
Opérateur de SMS & 1/2 & 3/10 & Non & Contrat type de l'opérateur \\
MEDISOFT (avenant 2026) & 2/2 & 10/10 & Oui & Avenant du 20~mai, audit
prévu en octobre \\
SAWA ASSURANCES (conjoint) & sans objet & sans objet & sans objet &
Convention non rédigée \\
\end{tableaucours}

\medskip
\textbf{Constatations.} Sur quatre contrats de sous-traitance, un seul
contient les douze clauses (MEDISOFT, depuis l'avenant de mai). SOREC
n'a aucune clause. L'éditeur du logiciel bancaire n'accepte pas d'audit et
ne mentionne pas le pays d'hébergement. Aucune convention de
responsabilité conjointe n'existe avec SAWA ASSURANCES.

\textbf{Analyse} (confiance élevée sur les contrats, faible sur
l'exécution). Les écarts contractuels sont établis. Pour l'exécution, seul
SOREC a été vu~: la conservation de cinq ans de classeurs est un écart à
l'article~13, dont NYANGA répond aussi (article~20).
\end{boxscenario}

\begin{boxlecture}{Le domaine D4, sous-traitants, à NYANGA}
\begin{tableaucours}{L{7.8cm}cL{3.2cm}}{%
    \tableauHeader{Exigence} &
    \tableauHeader{Cote} &
    \tableauHeader{Preuve}
}
Chaque sous-traitant est lié par un contrat & \OUI{} & Trois contrats \\
Les contrats définissent données et personnes (art.~30(3)) & \NON{} &
SOREC, opérateur \\
Les contrats contiennent les douze clauses & \NON{} & Un sur trois au
mieux \\
Le droit d'audit est prévu & \NON{} & Aucun contrat \\
SOREC détruit les données des dossiers clos & \NON{} & Classeurs
conservés \\
L'éditeur applique les mesures de sécurité convenues & \NEK{} & Pas
d'audit, pas de rapport \\
L'opérateur de SMS limite l'usage des numéros & \NEK{} & Aucune pièce \\
La convention avec SAWA ASSURANCES est publique & \NON{} & Inexistante \\
\end{tableaucours}

\medskip
\textbf{Décompte.} Un \OUI{}, cinq \NON{}, deux \NEK{} sur huit.
\hfill\lecture{0{,}125}{0{,}25}{1{,}30}

\medskip
\textbf{Lecture.} Domaine opaque ($U = 0{,}25$)~: on ne sait pas ce que font
deux des trois prestataires. Le manque de droit d'audit explique les deux
\NEK{}~: corriger les contrats est la condition pour lever l'incertitude.
\end{boxlecture}

% ============================================================================
\section{Ce que la Sous-traitance demande à chaque Métier}
\label{sec:c10-metiers}
% ============================================================================

\begin{boxmetier}{Juriste conformité}
Rédiger un avenant type aux douze clauses, le négocier en commençant par les
prestataires les plus exposés, et tenir à jour la liste des sous-traitants
avec leur contrat.
\end{boxmetier}

\begin{boxmetier}{Responsable de traitement}
Refuser ou quitter un prestataire qui refuse le droit d'audit pour un
traitement sensible. La responsabilité solidaire et l'article~20 font que
ses manquements sont aussi les vôtres.
\end{boxmetier}

\begin{boxmetier}{Auditeur de la protection des données}
Distinguer dans le rapport les écarts contractuels (établis sur pièces) et
les écarts d'exécution (établis seulement là où l'on a pu vérifier), et
coter \NEK{} l'exécution des prestataires non vus.
\end{boxmetier}

\begin{boxresume}
\begin{itemize}
    \item Tout recours à un sous-traitant exige un contrat qui définit au
          moins les données, les personnes, les obligations et les droits
          des parties (article~30).
    \item Le responsable s'assure des garanties du sous-traitant et veille à
          leur application (articles~16 et~30)~; les deux ont les mêmes
          obligations (article~20).
    \item Le cours déclare un critère de douze clauses, dont le droit
          d'audit, condition de toutes les autres.
    \item Les responsables conjoints fixent leurs obligations dans une
          convention rendue publique (article~31).
    \item À NYANGA, aucun contrat ne prévoit d'audit, et SOREC conserve cinq
          ans de données de débiteurs.
\end{itemize}
\end{boxresume}

\begin{boxinfo}
\textbf{Points de contrôle.} Vous avez compris ce chapitre si vous savez~:
\begin{itemize}
    \item contrôler un contrat contre les douze clauses~;
    \item distinguer sous-traitant et responsable conjoint~;
    \item choisir le cercle de vérification adapté au risque d'un
          prestataire.
\end{itemize}
\end{boxinfo}

\begin{boxexo}[title={\ding{45}\quad Exercice~13.1. Sous-traitant ou conjoint \niveaubase}]
Qualifiez la relation~: (a) un cabinet comptable qui établit la paie de
MBOLO~; (b) un laboratoire d'analyses qui reçoit les prélèvements et
envoie les résultats au médecin~; (c) l'hébergeur de NOUN PAY~; (d) NYANGA
et SAWA ASSURANCES pour l'assurance des emprunteurs.
\end{boxexo}

\begin{boxexo}[title={\ding{45}\quad Exercice~13.2. Lire un contrat \niveaumoyen}]
Le contrat de l'opérateur de SMS prévoit que « l'opérateur pourra utiliser
les données d'acheminement pour améliorer ses services et proposer des
offres à ses abonnés ». Qualifiez l'opérateur pour cette utilisation, et
dites quelles clauses vous ajoutez.
\end{boxexo}

\begin{boxexo}[title={\ding{45}\quad Exercice~13.3. La convention publique \niveaumoyen}]
Rédigez les rubriques de la convention entre NYANGA et SAWA ASSURANCES
exigée par l'article~31. Qui recueille le consentement au questionnaire de
santé, et sur quel fondement~?
\end{boxexo}

\begin{boxexo}[title={\ding{45}\quad Exercice~13.4. Le refus d'audit \niveauexpert}]
L'éditeur du logiciel bancaire refuse tout audit, invoquant la sécurité de
ses autres clients. Proposez trois solutions qui permettraient à NYANGA de
s'assurer des garanties exigées par l'article~30, et dites comment vous
coterez la proposition « l'éditeur applique les mesures convenues » dans
chaque cas.
\end{boxexo}

\bigskip

\begin{boxinfo}
\textbf{Pour aller plus loin.}
\begin{itemize}
    \item L'ISO/IEC~27701~\cite{iso27701}, pour les obligations des
          sous-traitants.
    \item L'ISO/IEC~27018~\cite{iso27018}, pour l'hébergement en nuage.
\end{itemize}
\end{boxinfo}

\begin{boxinfo}
\textbf{Transition.} L'éditeur du logiciel bancaire héberge les comptes à
Abidjan, et l'hébergeur de NOUN PAY en Europe. Les données quittent le
Cameroun~: c'est un transfert, soumis à autorisation. C'est l'objet du
Chapitre~\ref{chap:transferts}.
\end{boxinfo}
